\skyrius{Produkto reikalavimai ir projektavimas}
\subsection{Reikalavimai ir sėkmės kriterijai}
\gaires{Prieš kurdami nuspręskite, kaip atpažinsite tinkamą rezultatą.
Funkcinis reikalavimas nusako, ką produktas daro, o kokybės reikalavimas -- kaip
gerai ar kokiomis sąlygomis tai daro. Kriterijus „patogu“ pakeiskite tikrinamu
požymiu. R1, R2 ir kiti žymenys leis susieti reikalavimus su bandymais.}
\begin{table}[H]
\caption{Produkto reikalavimai ir jų patikros planas}
\label{tab:reikalavimai}
\begin{tabularx}{\textwidth}{@{}p{0.8cm}YYY@{}}
\toprule
ID & Reikalavimas ir jo pagrindas & Priėmimo kriterijus & Patikros būdas \\
\midrule
R1 & \pildyti{Funkcija ir poreikis} & \pildyti{Stebimas rezultatas} & \pildyti{Bandymas T1} \\
R2 & \pildyti{Kokybės reikalavimas} & \pildyti{Pamatuojama riba} & \pildyti{Bandymas T2} \\
\bottomrule
\end{tabularx}
\end{table}

\subsection{Sandara ir technologijų pasirinkimas}
\pildyti{Aprašykite produkto dalis, jų ryšius ir veikimo principą. Informacinei
sistemai tinka architektūros, duomenų ar sąsajos schema; fiziniam produktui --
brėžinys, komponentų ir medžiagų pagrindimas. Palyginkite realias alternatyvas
ir paaiškinkite, kodėl pasirinkimas atitinka reikalavimus.}
\begin{figure}[H]
\centering
\fbox{\parbox[c][2.2cm][c]{0.86\textwidth}{\centering
\pildyti{Įkelkite savo architektūros, duomenų, sąsajos ar konstrukcijos schemą.}}}
% Rėmelį keiskite, pvz.: \includegraphics[width=0.86\textwidth]{media/schema.pdf}
\caption{Projektuojamo produkto sandara}
\label{fig:sandara}
\saltinis{\pildyti{autoriaus schema arba naudoto šaltinio nuoroda}}
\end{figure}
\pildyti{Paaiškinkite \ref{fig:sandara} paveikslo dalis ir svarbiausius ryšius.}

\subsection{Darbo planas ir rizikos}
\pildyti{Trumpai apibūdinkite kūrimo etapus, numatytus tarpinius rezultatus ir
rizikas. Detalus planas pateikiamas \ref{app:planas} priede. Paaiškinkite,
kaip spręsite aktualiausias rizikas: laiko, duomenų, įrangos ar technologijų.}
